\documentclass[11pt]{article}
\usepackage[margin=1in]{geometry}
\usepackage{amsmath,amssymb,amsthm}
\usepackage{graphicx}
\usepackage[hidelinks]{hyperref}

\newtheorem{theorem}{Theorem}
\newtheorem{lemma}[theorem]{Lemma}

\newcommand{\R}{\mathbb{R}}
\newcommand{\D}{\mathcal{D}}
\newcommand{\Hsp}{\mathcal{H}_{p'}^{t',r'}}
\newcommand{\HY}{\mathcal{H}_{p'}^{t',r'}(Y)}
\newcommand{\Hs}{\mathcal{H}_{p'}^{t',r'}}
\newcommand{\norm}[1]{\left\|#1\right\|}

\begin{document}

\begin{center}
{\Large Full positive answer to Problem 4.9 of arXiv:2505.19143}\\[0.5em]
{\large Vector-valued Bourgain-Morrey block spaces are norm-solid}\\[0.5em]
Agent: \texttt{agent\_lane\_15}; run: \texttt{fa\_banach\_001}
\end{center}

\section*{Source and Question}

The source paper is Tengfei Bai, Pengfei Guo, and Jingshi Xu,
\emph{The preduals of Banach space valued Bourgain-Morrey spaces},
arXiv:2505.19143 \cite{BGX}.

Problem 4.9 asks whether the vector-valued block space has the lattice
property under the pointwise norm order. In the notation of the source, let
\(X\) be a Banach space with a predual \({}^\ast X\), let \({}^\ast X\) be
reflexive or \(X\) be separable, and let
\[
1<p<t<r<\infty \quad\text{or}\quad 1<p\leq t<r=\infty .
\]
If \(g\in \Hsp({}^\ast X)\) and \(f\) is a \({}^\ast X\)-valued function with
\[
\norm{f(x)}_{{}^\ast X}\leq \norm{g(x)}_{{}^\ast X}
\quad\text{for a.e. }x,
\]
must \(f\) also belong to \(\Hsp({}^\ast X)\)?

\begin{center}
\includegraphics[width=0.96\textwidth]{figures/open_problem_crop.png}
\end{center}

\section*{Answer}

Yes. In fact, the vector-valued block norm is exactly the scalar block norm of
the pointwise Banach norm:
\[
\norm{F}_{\Hsp(Y)}=\norm{\,\norm{F(\cdot)}_Y\,}_{\Hsp}
\]
for every Banach space \(Y\) and every strongly measurable \(Y\)-valued
function \(F\) for which either side is finite. Consequently the desired
lattice property holds for arbitrary Banach-valued block spaces; no coordinate
lattice or order structure on \(Y\) is needed.

\section*{Proof Intuition}

The source's block definition measures a vector-valued block only through the
scalar function \(x\mapsto \norm{b(x)}\). Therefore directions in the Banach
space should not matter. One direction is obtained by taking norms in any
vector block decomposition. The reverse direction lifts a scalar decomposition
of \(\norm{F(\cdot)}\) by multiplying each scalar block by the measurable unit
direction \(F(x)/\norm{F(x)}\). Solidity then follows by cutting the scalar
positive decomposition of \(\norm{g(\cdot)}\) down to \(\norm{f(\cdot)}\).

\section*{Definitions}

We use the source's dyadic block notation. A scalar \((p',t')\)-block is a
measurable scalar function \(a\), supported in a dyadic cube \(Q\), such that
\[
\norm{a}_{L^{p'}}\leq |Q|^{1/t-1/p}.
\]
For a Banach space \(Y\), a \(Y\)-valued \((p',t',Y)\)-block is a strongly
measurable \(Y\)-valued function \(b\), supported in a dyadic cube \(Q\), such
that
\[
\norm{\,\norm{b(\cdot)}_Y\,}_{L^{p'}}\leq |Q|^{1/t-1/p}.
\]
The block space \(\Hsp(Y)\) consists of those \(Y\)-valued functions admitting
an almost everywhere convergent representation
\[
F=\sum_{(j,k)\in \mathbb Z^{n+1}}\lambda_{j,k} b_{j,k},
\qquad \{\lambda_{j,k}\}\in \ell^{r'},
\]
where \(b_{j,k}\) is a \((p',t',Y)\)-block supported on \(Q_{j,k}\). Its norm is
the infimum of \(\norm{\lambda}_{\ell^{r'}}\) over all such representations.
The scalar space \(\Hsp\) is the case \(Y=\R\).

\section*{Scalar Solidity}

\begin{lemma}[Scalar solidity]
Let \(0\leq u\leq v\) almost everywhere and \(v\in\Hsp\). Then
\[
u\in\Hsp,\qquad \norm{u}_{\Hsp}\leq \norm{v}_{\Hsp}.
\]
\end{lemma}

\begin{proof}
Fix \(\varepsilon>0\). Choose an admissible scalar block representation
\[
v=\sum_i \lambda_i a_i,\qquad
\norm{(\lambda_i)}_{\ell^{r'}}\leq \norm{v}_{\Hsp}+\varepsilon .
\]
Here the index \(i\) abbreviates \((j,k)\). Put
\[
A(x)=\sum_i |\lambda_i|\, |a_i(x)|.
\]
By the same block-summability argument used in the source for admissible block
sums, \(A\) is finite almost everywhere and belongs to \(\Hsp\) with norm at
most \(\norm{(\lambda_i)}_{\ell^{r'}}\). Since \(v\geq0\) and
\(v=\sum_i\lambda_i a_i\), the triangle inequality gives the pointwise
majorization \(0\leq u\leq v=|v|\leq A\).

Define
\[
c_i(x)=
\begin{cases}
\dfrac{u(x)}{A(x)}\, |a_i(x)|, & A(x)>0,\\[0.8em]
0, & A(x)=0.
\end{cases}
\]
Then \(0\leq c_i\leq |a_i|\), so each \(c_i\) is a scalar block supported on
the same cube as \(a_i\). Also,
\[
u(x)=\sum_i |\lambda_i| c_i(x)
\]
for almost every \(x\). Hence \(u\) has an admissible block representation with
coefficient sequence \((|\lambda_i|)\), and
\[
\norm{u}_{\Hsp}\leq \norm{(\lambda_i)}_{\ell^{r'}}
\leq \norm{v}_{\Hsp}+\varepsilon .
\]
Letting \(\varepsilon\downarrow0\) proves the claim.
\end{proof}

\section*{Norm Reduction}

\begin{lemma}[Pointwise norm reduction]
Let \(Y\) be a Banach space and let \(F\) be a strongly measurable
\(Y\)-valued function. Then
\[
F\in\Hsp(Y)
\quad\Longleftrightarrow\quad
\norm{F(\cdot)}_Y\in\Hsp,
\]
and in this case
\[
\norm{F}_{\Hsp(Y)}=\norm{\,\norm{F(\cdot)}_Y\,}_{\Hsp}.
\]
\end{lemma}

\begin{proof}
First suppose \(F\in\Hsp(Y)\). Fix \(\varepsilon>0\) and choose
\[
F=\sum_i \lambda_i b_i,\qquad
\norm{(\lambda_i)}_{\ell^{r'}}\leq \norm{F}_{\Hsp(Y)}+\varepsilon,
\]
with \(Y\)-valued blocks \(b_i\). Set
\[
A(x)=\sum_i |\lambda_i|\,\norm{b_i(x)}_Y .
\]
The scalar functions \(\norm{b_i(\cdot)}_Y\) are scalar blocks, so \(A\in\Hsp\)
with
\[
\norm{A}_{\Hsp}\leq \norm{(\lambda_i)}_{\ell^{r'}}.
\]
Moreover \(\norm{F(x)}_Y\leq A(x)\) almost everywhere. By scalar solidity,
\[
\norm{\,\norm{F(\cdot)}_Y\,}_{\Hsp}
\leq \norm{(\lambda_i)}_{\ell^{r'}}
\leq \norm{F}_{\Hsp(Y)}+\varepsilon .
\]
Letting \(\varepsilon\downarrow0\) gives
\[
\norm{\,\norm{F(\cdot)}_Y\,}_{\Hsp}\leq \norm{F}_{\Hsp(Y)}.
\]

Conversely assume \(\phi=\norm{F(\cdot)}_Y\in\Hsp\). Define the strongly
measurable direction
\[
u(x)=
\begin{cases}
F(x)/\phi(x), & \phi(x)>0,\\
0, & \phi(x)=0 .
\end{cases}
\]
Then \(\norm{u(x)}_Y\leq1\) and \(F(x)=\phi(x)u(x)\) almost everywhere. Choose
an admissible scalar representation
\[
\phi=\sum_i \mu_i a_i,\qquad
\norm{(\mu_i)}_{\ell^{r'}}\leq \norm{\phi}_{\Hsp}+\varepsilon.
\]
For each \(i\), put \(B_i(x)=a_i(x)u(x)\). Since
\[
\norm{B_i(x)}_Y\leq |a_i(x)|,
\]
each \(B_i\) is a \(Y\)-valued block supported on the same cube as \(a_i\).
The scalar convergence of \(\sum_i\mu_i a_i\) to \(\phi\) implies convergence
of \(\sum_i\mu_i B_i\) to \(\phi u=F\) in \(Y\) almost everywhere. Therefore
\[
\norm{F}_{\Hsp(Y)}\leq \norm{(\mu_i)}_{\ell^{r'}}
\leq \norm{\phi}_{\Hsp}+\varepsilon .
\]
Letting \(\varepsilon\downarrow0\) gives the reverse inequality and hence the
norm identity.
\end{proof}

\section*{Main Theorem}

\begin{theorem}[Affirmative answer to Problem 4.9]
Let \(X\) be a Banach space such that \(({}^\ast X)^\ast=X\). Assume
\({}^\ast X\) is reflexive or \(X\) is separable, and let
\[
1<p<t<r<\infty \quad\text{or}\quad 1<p\leq t<r=\infty .
\]
If \(g\in \Hsp({}^\ast X)\) and \(f\) is a \({}^\ast X\)-valued measurable
function satisfying
\[
\norm{f(x)}_{{}^\ast X}\leq\norm{g(x)}_{{}^\ast X}
\quad\text{for almost every }x\in\R^n,
\]
then \(f\in \Hsp({}^\ast X)\). More precisely,
\[
\norm{f}_{\Hsp({}^\ast X)}
\leq
\norm{g}_{\Hsp({}^\ast X)}.
\]
\end{theorem}

\begin{proof}
Apply the pointwise norm reduction lemma to \(g\). It gives
\[
\norm{g(\cdot)}_{{}^\ast X}\in\Hsp
\quad\text{and}\quad
\norm{\,\norm{g(\cdot)}_{{}^\ast X}\,}_{\Hsp}
=\norm{g}_{\Hsp({}^\ast X)} .
\]
Since
\[
0\leq \norm{f(\cdot)}_{{}^\ast X}
\leq \norm{g(\cdot)}_{{}^\ast X}
\]
almost everywhere, scalar solidity gives
\[
\norm{f(\cdot)}_{{}^\ast X}\in\Hsp,\qquad
\norm{\,\norm{f(\cdot)}_{{}^\ast X}\,}_{\Hsp}
\leq
\norm{\,\norm{g(\cdot)}_{{}^\ast X}\,}_{\Hsp}.
\]
Applying the pointwise norm reduction lemma again, now to \(f\), yields
\[
f\in\Hsp({}^\ast X)
\]
and the asserted norm inequality.
\end{proof}

\section*{Verification Notes and Limitations}

The proof uses only the block definition from the source paper and does not
use an order or coordinate lattice on \({}^\ast X\). Thus it proves the source
problem in the stated range and in fact applies to arbitrary Banach targets
\(Y\), whenever the corresponding block space is defined with the same
pointwise-norm block condition.

No computational verification is needed. The main technical points for human
review are:
\begin{itemize}
\item the scalar cut-down construction in the scalar solidity lemma;
\item measurability of the direction \(F/\norm{F}\), which follows from strong
      measurability of \(F\);
\item preservation of block bounds under \(a_i\mapsto a_i u\).
\end{itemize}

\section*{Novelty Check}

Before promotion, the run indexes were searched for arXiv:2505.19143, the
paper title, Bourgain-Morrey block spaces, and lattice-property keywords. No
prior packet or attempt was found. The local parsed arXiv corpus was searched
for exact phrases from Problem 4.9, including vector-valued Bourgain-Morrey
block-space lattice-property phrases; only the source paper surfaced as a
direct hit. Web searches on 2026-06-28 for exact phrases including
``Problem 4.9 Bourgain-Morrey lattice property'' and
``Banach space valued Bourgain-Morrey lattice property'' returned the source
arXiv record and no separate answer in this bounded pass.

\begin{thebibliography}{9}

\bibitem{BGX}
T. Bai, P. Guo, and J. Xu,
\emph{The preduals of Banach space valued Bourgain-Morrey spaces},
arXiv:2505.19143, 2025.

\end{thebibliography}

\end{document}
